%% Chapter 2: Literature Review: chapters/chapter02.tex

% ==============================================================================
% CONTENT BLUEPRINT — Chapter 2: Literature Review (M.S. Thesis)
% ==============================================================================
% Reading audience: examiner who needs to assess whether the candidate has
% read broadly and positioned the work correctly.
% Goal: deliver a critical synthesis of prior work — not a list of abstracts.
%
% Recommended length: 12–20 pages.
%
% ----------------------------------------------------------------------------
% §2.1 Historical Perspective  (2–3 pages)
%   • Trace the field from its earliest principled formulation. Include:
%       – The breakthrough paper(s) that made the field quantitative.
%       – Major theoretical milestones (no need to recite all of them; pick
%         the ones whose equations §3.1 will reuse).
%       – A one-paragraph "why this line of work eventually stalled or
%         diverged" closing the section.
%
% §2.2 Theoretical Foundations  (3–4 pages)
%   • Develop the formal framework that chapter 4 will need:
%       – State the central equations (numbered, with units in \si{...}).
%       – For each, derive the working formula that your code implements.
%       – Identify the regimes of validity (e.g. "valid for ε ≤ 5 % biaxial
%         strain; above this, anharmonic contributions dominate").
%   • End with a small worked example, ready for §4.2 to compare against.
%
% §2.3 Computational and Experimental Methods  (3–4 pages)
%   • Comprehensive review of the tools other groups have used. For each:
%       – Method name, who introduced it, what it computes.
%       – Strengths and weaknesses in 1–2 sentences.
%       – At least one citation per mainstream method (DFT, GW, BSE, MD,
%         Monte Carlo, etc. — only those relevant to your work).
%   • Identify the "method gap" if any (e.g. "few works combine hybrid
%     functionals with spin-orbit coupling for this material class").
%
% §2.4 Material-Specific Body of Work  (3–4 pages)
%   • Critical summary of the closest-studied systems. One subsection per
%     relevant sub-class (e.g. "2.4.1 TMD monolayers", "2.4.2 Heterostructures").
%   • For each subsection:
%       – A representative comparison table (literature | method | reported
%         value | reference).
%       – A discussion paragraph that identifies what is consistent vs what
%         is still debated.
%       – Forward-pointer to what the present thesis adds.
%
% §2.5 Positioning of the Present Work  (½–1 page)
%   • One paragraph that names exactly what is new relative to §2.1–§2.4.
%   • Cite this section in §1.5 / §5.2 of the introduction/conclusion to
%     keep the narrative thread tight.
%
% Writing-style rules:
%   • Critical, not descriptive. "X used DFT with LDA and obtained bandgap
%     E_g = 0.6 eV; however, LDA systematically underestimates gaps — the
%     value should be re-examined using hybrid functionals (see §2.3)."
%   • Compare studies against each other, not just one against an ideal.
%   • Cite 40–80 references across the chapter (typical M.S. range).
% ==============================================================================

\chapter{Literature Review}
\label{ch:lit_review}

\section{Historical Perspective}
The theoretical foundations of electronic structure calculations date back to the pioneering works of Hohenberg, Kohn, and Sham in the mid-1960s \citep{hohenberg1964, kohn1965}. Their development of \gls{dft} revolutionized quantum chemistry and solid-state physics by representing the complex many-body Schrödinger equation in terms of single-particle equations based on electron density.

\section{Developments in Semiconductor Physics}
Early studies of semiconductor properties relied heavily on empirical pseudopotential methods. Over the last two decades, with the increase of computing power, first-principles methods have become the standard. For example, \citet{kresse1996} developed the projector augmented-wave method implemented in \gls{vasp}, enabling high-precision simulations of complex materials.

\section{Strain Engineering in Two-Dimensional Materials}
Strain engineering has emerged as a promising avenue to tune the properties of low-dimensional semiconductors. Biaxial and uniaxial strain have been shown to induce direct-to-indirect bandgap transitions in transition metal dichalcogenides \citep{novoselov2005}. However, most studies have focused on simple structures. The behavior of multi-component superlattices under high strain fields remains an active area of investigation. This work aims to address this gap by simulating realistic superlattice structures under realistic strain profiles.
